% !PRINTED 773
\weberpage{773}
%<<CONT
then the $c_1, c_2, \ldots$ are \gesperrt{integers}, because $\chi(x)$ must be a divisor of $x^m - 1$ (by \gesperrt{Gauss}'s theorem, Bd.~I, \S.\,2).

From this it follows that, if $\varphi(x)$ denotes an integral function of $x$ with integral rational coefficients, the coefficients $m_0, m_1, \ldots, m_{\nu-1}$ in (9) become integers, even when the degree of $\varphi(x)$ is originally higher than $\nu$; for one obtains these coefficients by seeking out the residue of the division of $\varphi(x)$ by $\chi(x)$.

We now assume that $f_n(x)$ is decomposable in the field $R(\alpha)$, and that accordingly:

\weq{11}{f_n(x) = \varphi(x,\alpha)\ \psi(x,\alpha),}

where $\varphi(x,\alpha)$ and $\psi(x,\alpha)$ are two integral functions of $x$, of which neither is constant, and whose coefficients are contained in $R(\alpha)$. From (8) and (11) it follows, if we set $x=1$,

\weq{12}{p = \varphi(1,\alpha)\ \psi(1,\alpha),}

where $\varphi(1,\alpha)$, $\psi(1,\alpha)$ can be represented, as numbers of the field $R(\alpha)$, in the form (9). If we denote the least common denominators of these two representations by $a$ and $b$, we obtain:

\weq{13}{\begin{aligned}
a\,\varphi(1,\alpha) &= a_0 + a_1\alpha + a_2\alpha^2 + \cdots + a_{\nu-1}\alpha^{\nu-1} = A(\alpha)\\
b\,\psi(1,\alpha) &= b_0 + b_1\alpha + b_2\alpha^2 + \cdots + b_{\nu-1}\alpha^{\nu-1} = B(\alpha),
\end{aligned}}

where $a, a_0, a_1, \ldots, a_{\nu-1}$ are integers without a common divisor and $A$ is a function symbol. The same holds for $b, b_0, b_1, \ldots, b_{\nu-1}$ and $B$.

The roots of the equation $\varphi(x,\alpha) = 0$ are to be found among the primitive $n$th roots of unity, and these are all powers of one among them, $r$. Accordingly there are certain powers $r^h, r^{h'}, r^{h''}, \ldots$ such that

\[
\varphi(x,\alpha) = (x - r^h)(x - r^{h'})(x - r^{h''}) \ldots
\]

holds, and consequently, by (13):

\[
A(\alpha) = a(1 - r^h)(1 - r^{h'})(1 - r^{h''}) \ldots
\]

If we take the $n$th power of this, it follows by (1), having regard to the equation $r^n = 1$:

\[
[A(\alpha)]^n = p\,\chi(r),
\]

where $\chi(x)$ denotes an integral function of the variable $x$ with integral coefficients.
